% GENERATED FILE. Edit canonical lessons, then run npm run generate:books.
% canonical-source-hash: fnv1a64:4b7de9381ea7188f
% canonical-lessons: FR-C261-je-devrais, FR-C261-expliquer, FR-C261-revision-explication

\chapter{Should, and Why — je devrais, expliquer}
\label{ch:fr-should-and-why}
\hlchaptermodality{261}

\begin{chapteropening}
\textbf{By the end of this chapter:} \emph{I can give advice with je devrais and tu devrais, and explain a decision in four moves.}

Complete FR-C261-revision-explication: a whole explanation aloud: a view, a cause, a consequence and a piece of advice.
\end{chapteropening}

\section[je devrais]{je devrais — what you ought to do about it}

\label{lesson:FR-C261-je-devrais}



\begin{quote}
Say \textquotedblleft{}I must leave\textquotedblright{}. (\emph{Je dois partir.}) Then ask a stranger for a coffee. (\emph{Je voudrais un café.})

An explanation often ends on a third move: what somebody ought to do. That move has its own form, and you already own its ending.
\end{quote}

\subsection*{You'll want to know: je devrais}
\textbf{je devrais} — \textquotedblleft{}I should\textquotedblright{}.

It is \emph{devoir} with the soft ending of \emph{je voudrais}: \textbf{je devrais}, \textbf{tu devrais}, \textbf{il devrait}, \textbf{vous devriez}. A bare infinitive follows, with nothing in between, exactly as after \emph{je dois}:

\begin{quote}
\textbf{Je devrais} travailler. — I should work. \textbf{Tu devrais} dormir. — You should sleep. \emph{(to a friend)} \textbf{Vous devriez} vous reposer. — You should rest. \emph{(to anyone you call vous)}
\end{quote}

\emph{Je dois travailler} has no way out. \emph{Je devrais travailler} says I know I ought to, and have not promised. That softness is the register advice usually wants. The ending is the one in \emph{voudrais}, and chapter two hundred and sixty-seven gives the whole form its name and its recipe.

\begin{grammarlens}[title={Grammar Lens: advice, with its reason}]
Advice sits naturally at the end of an arrow:

\begin{quote}
\emph{Tu es fatigué, \textbf{donc tu devrais} dormir.} — You're tired, so you should sleep. \emph{\textbf{Je pense que tu devrais} te reposer.} — I think you should rest.
\end{quote}

With \textbf{on}, it becomes \textquotedblleft{}we ought to\textquotedblright{}: \emph{\textbf{On devrait} partir.} — We ought to leave. Said late in an evening, it is a polite way to end it.
\end{grammarlens}

\begin{cousinweb}
\textbf{devoir} is inherited from Latin \emph{dēbēre}, \textquotedblleft{}to owe\textquotedblright{} (EtymDB 2.1), and \emph{dēbēre} is a compound: \emph{dē}, \textquotedblleft{}from\textquotedblright{}, plus \emph{habēre}, \textquotedblleft{}to have\textquotedblright{} (Lewis \& Short, under \emph{debeo}). To owe is to \textbf{have something from someone}: you are holding what is theirs. So \emph{je devrais} is, at root, \textquotedblleft{}I would owe it\textquotedblright{}: the debt is there, and nobody is collecting yet.
\end{cousinweb}

\subsection*{Your turn}
Say these aloud:

\begin{itemize}
\raggedright
  \item \emph{je devrais, tu devrais, il devrait}
  \item the hard one, then the soft one — \emph{Je dois travailler.} then \emph{Je devrais travailler.}
  \item advice to a tired friend — \emph{Tu es fatigué, alors tu devrais dormir.}
  \item \emph{On devrait partir.}
\end{itemize}

\begin{tcolorbox}[breakable,colback=teal!4,colframe=teal!35!black,title={Before you move on}]
What is the difference between \emph{je dois} and \emph{je devrais}? (\emph{\textbf{Je dois}} is \textquotedblleft{}I must\textquotedblright{}; \emph{\textbf{je devrais}} is \textquotedblleft{}I should\textquotedblright{}, softer.) Which form you know has the same ending? (\emph{\textbf{Je voudrais.}}) What are the two Latin pieces of \emph{dēbēre}? (\emph{\textbf{Dē}} and \emph{\textbf{habēre}} — to have from someone.)
\end{tcolorbox}
\section[expliquer pourquoi]{expliquer pourquoi — to unfold}

\label{lesson:FR-C261-expliquer}



\begin{quote}
Four pieces are in your hands: \emph{je pense que} for a view, \emph{parce que} and \emph{donc} for the two arrows, \emph{je devrais} for advice. This lesson puts them in a row, under the verb that names the act.
\end{quote}

\begin{grammarlens}[title={Grammar Lens: explaining something to someone}]
You have \textbf{expliquer}, \textquotedblleft{}to explain\textquotedblright{}, a regular \emph{-er} verb. The thing explained comes straight after it; the person who hears it takes \textbf{à}, or stands in front as a small pronoun:

\begin{quote}
\emph{J'explique le problème \textbf{à} ma mère.} — I explain the problem to my mother. \emph{Je \textbf{t'}explique.} — Let me explain. \emph{(to a friend)} \emph{Vous pouvez \textbf{m'}expliquer ?} — Could you explain it to me?
\end{quote}

And the explanation itself often opens with \textbf{la raison}, \textquotedblleft{}the reason\textquotedblright{}: \emph{La raison, c'est que…} — \textquotedblleft{}The reason is that…\textquotedblright{}
\end{grammarlens}

\begin{grammarlens}[title={Grammar Lens: the shape of an explanation}]
Here is a whole explanation, every word of it met before:

\begin{quote}
— Pourquoi tu ne sors pas ? — \textbf{Je pense qu'}il va pleuvoir. Et je suis malade, \textbf{parce qu'}hier il faisait froid. \textbf{Donc} je reste à la maison. \textbf{Je devrais} dormir.
\end{quote}

Four moves, and a word for each: a \textbf{view} (\emph{je pense que}), a \textbf{cause} pointing back (\emph{parce que}), a \textbf{consequence} pointing forward (\emph{donc}), and \textbf{advice} (\emph{je devrais}). With both arrows you can build it cause first or claim first, and the listener can follow either.
\end{grammarlens}

\begin{cousinweb}
\textbf{expliquer} is a learned borrowing from Latin \emph{explicāre}, \textquotedblleft{}to unfold, to open out flat\textquotedblright{} (EtymDB 2.1; Lewis \& Short, under \emph{explico}), which is \emph{ex-}, \textquotedblleft{}out\textquotedblright{}, plus \emph{plicāre}, \textquotedblleft{}to fold\textquotedblright{}. \textbf{To explain is to unfold}: something was folded shut, and you open it out.

The same \emph{plicāre} came down into French the slow way, by mouth, and gave \textbf{plier}, \textquotedblleft{}to fold\textquotedblright{}, which you have (EtymDB 2.1). English kept the Latin in \textbf{explicit}, and \textbf{complicated} is the opposite: folded together.
\end{cousinweb}

\subsection*{Your turn}
Say these aloud:

\begin{itemize}
\raggedright
  \item \emph{Je t'explique.} — then \emph{Vous pouvez m'expliquer ?}
  \item the four moves — \emph{Je pense qu'il va pleuvoir. Je suis malade. Donc je reste à la maison. Je devrais dormir.}
  \item the same, cause last — \emph{Je reste à la maison parce que je suis malade.}
  \item \emph{La raison, c'est que je suis malade.}
\end{itemize}

\begin{tcolorbox}[breakable,colback=teal!4,colframe=teal!35!black,title={Before you move on}]
Name the four moves of an explanation and a word for each. (View \emph{\textbf{je pense que}}, cause \emph{\textbf{parce que}}, consequence \emph{\textbf{donc}}, advice \emph{\textbf{je devrais}}.) Who takes \emph{à} after \emph{expliquer}? (The \textbf{person} who hears it.) What is \emph{expliquer}, taken apart? (\emph{\textbf{Ex-}} and \emph{\textbf{plicāre}}: \textbf{to unfold}.)
\end{tcolorbox}
\section[(review)]{Review — unfolding a decision}

\label{lesson:FR-C261-revision-explication}



\begin{quote}
Say \textquotedblleft{}I should work\textquotedblright{}. (\emph{Je devrais travailler.}) Then \textquotedblleft{}let me explain\textquotedblright{}, to a friend. (\emph{Je t'explique.})
\end{quote}

\begin{grammarlens}[title={What you've built}]
Two chapters, and a whole explanation. A view: \emph{je pense que}, \emph{à mon avis}. A cause, pointing back: \emph{parce que}. A consequence, pointing forward: \emph{donc}, \emph{alors}. And advice, softened: \emph{je devrais}, \emph{tu devrais}, \emph{on devrait}.

Here is a friend's question, and an answer in four moves:

\begin{quote}
— \emph{Pourquoi tu ne viens pas au cinéma ?} — \emph{À mon avis, le film est ennuyeux. Et je suis fatigué, parce que j'ai beaucoup travaillé. Donc je reste à la maison. Mais toi, tu devrais y aller !}
\end{quote}

\emph{Tu ne viens pas} is from \textbf{venir}, \textquotedblleft{}to come\textquotedblright{}, one of the verbs that take \emph{être} in the passé composé.
\end{grammarlens}

\subsection*{Your turn}
Say these aloud:

\begin{itemize}
\raggedright
  \item the view — \emph{À mon avis, le film est ennuyeux.}
  \item the cause — \emph{Je suis fatigué, parce que j'ai beaucoup travaillé.}
  \item the consequence — \emph{Donc je reste à la maison.}
  \item the advice — \emph{Mais toi, tu devrais y aller !}
  \item an explanation of your own, in four moves, for something you are not doing tonight
\end{itemize}

\begin{tcolorbox}[breakable,colback=teal!4,colframe=teal!35!black,title={Before you move on}]
Which move does \emph{je devrais} make? (\textbf{Advice.}) Say \textquotedblleft{}could you explain it to me?\textquotedblright{} (\emph{\textbf{Vous pouvez m'expliquer ?}}) What does \emph{expliquer} do, at root? (It \textbf{unfolds}.)
\end{tcolorbox}
